PhysAgent targets the physics specification stage---the SDevice \texttt{Physics} block and \texttt{.par} file---while structure (\texttt{.tdr}), solver, and bias schedule remain fixed, orthogonal to prior structure-generation work~\cite{nguyen, malts}. Three generators share the same base LLM and prompt: \textit{Vanilla} (spec only), \textit{LLM+Guide} (spec + Sentaurus manual), and \textit{PhysAgent} (spec + ontology + literature). We validate three claims.

\textbf{Successful TCAD Execution.}
On a Si nMOSFET we compare three modes (Vanilla, LLM+Guide, PhysAgent) across three prompt levels (L1--L3: explicit model list, categories, intent only), three trials each. Vanilla and LLM+Guide produce characteristic syntactic errors---invented model names, undefined parameters, malformed material blocks---lacking any structural anchor, whereas PhysAgent copies directly from typed ontology nodes at {\color{red}48$\times$} smaller prefill than LLM+Guide. Across L1$\to$L3 model coverage tracks prompt specificity: which physical mechanisms the deck selects depends on what the prompt asserts, making phenomenon-level prompting central to the simulation's analytical scope ({\color{red}Fig.~\ref{fig:exp1}}).

\textbf{Literature-Grounded Parameter File Generation.}

\textbf{Phenomenon-Driven Model Selection and Deck Construction.}
